\documentclass{article}
\usepackage[utf8]{inputenc}
\usepackage[margin=2cm,includefoot]{geometry}
\usepackage{booktabs}
\usepackage{graphicx}
\usepackage{amsmath, amssymb, amsthm}
\usepackage{physics}
\usepackage[shortlabels]{enumitem}
\usepackage{cancel}
\usepackage{amssymb}
\usepackage{parskip}
\usepackage{braket}

\let\oldemptyset\emptyset
\let\emptyset\varnothing

\usepackage{fancyhdr}
\pagestyle{fancy}
\lhead{UCSCIPHY25 - Statistical and quantum physics 2026}
\rhead{Adam Zich (0074187)}
\fancyfoot{}
\fancyfoot[R]{}

\usepackage{listings}
\usepackage{color}

\newcommand{\R}{\mathbb{R}}
\newcommand{\p}{\mathbb{P}}
\newcommand{\N}{\mathbb{N}}
\newcommand{\Z}{\mathbb{Z}}
\newcommand{\Q}{\mathbb{Q}}
\newcommand{\I}{\mathbb{I}}
\newcommand{\C}{\mathbb{C}}
\newcommand{\E}{\mathbb{E}}
\newcommand{\sP}{\mathcal{P}}
\newcommand{\mbeq}{\overset{!}{=}}

\lstset{frame=tb,
  language=java,
  aboveskip=3mm,
  belowskip=3mm,
  showstringspaces=false,
  columns=flexible,
  basicstyle={\small\ttfamily},
  numbers=none,
  numberstyle=\tiny\color{gray},
  keywordstyle=\color{blue},
  commentstyle=\color{dkgreen},
  stringstyle=\color{mauve},
  breaklines=true,
  breakatwhitespace=true,
  tabsize=3
}

\begin{document}
\begin{center}
 \LARGE{\bfseries{ Wednesday (Week 1 - 1) 11.2.2026 }}
 \line(1,0){430}
\end{center}

\section{Exercise set}
    \subsection{The mass of the atmosphere}
        Atmospheric pressure is $p_{air} \approx 100000Pa$. The radius of the Earth is $R = 6370000m$. We calculate the height of the atmosphere from the pressure as
        \begin{align*}
            p &= \rho hg,\\
            m &= \rho V,\\
            p &= \dfrac{m}{A}g,\\
            m &= \dfrac{pA}{g} = \dfrac{100000 \cdot 4 \cdot \pi \cdot 6378000^2}{9.81} = 5.21 \cdot 10^{18}kg.
        \end{align*}
        
    \subsection{Atmospheric pressure for non constant temperature}

    \subsection{Definition of flow lines}
    \begin{enumerate}[label=\alph*]
        \item 
    \end{enumerate}

    \section{Exercise set}
    \section{Exercise set}
    \section{Exercise set}
    \section{Exercise set}

        \subsection{Storm surges}

            \begin{enumerate}[label=\alph*)]
                \item Scale analysis
                    The continuity equation says that
                    \begin{align*}
                        \dfrac{\partial \rho}{\partial t} + \nabla \cdot (\rho \mathbf{u}) = 0.
                    \end{align*}
                    Since we have constant density and $\partial_yQ = 0, \forall Q$.

                    We can assume that
                    \begin{align*}
                        \frac{W}{U} \approx 0,
                    \end{align*}
                    since vertical velocities are typically smaller than horizontal ones.
                \item 
                \item 
            \end{enumerate}

\end{document}
